%%%%%%%%%%%%%%%%%%%%%%%%%%%%%%%%%%%%%%%%%%%%%%%%%%%%%%%%%%%%%%%%%%%%%%%%%%%%%%%%
%2345678901234567890123456789012345678901234567890123456789012345678901234567890
%        1         2         3         4         5         6         7         8

\documentclass[letterpaper, 10 pt, conference]{ieeeconf}  % Comment this line out if you need a4paper

%\documentclass[a4paper, 10pt, conference]{ieeeconf}      % Use this line for a4 paper

\IEEEoverridecommandlockouts                              % This command is only needed if 
                                                          % you want to use the \thanks command

\overrideIEEEmargins                                      % Needed to meet printer requirements.

%In case you encounter the following error:
%Error 1010 The PDF file may be corrupt (unable to open PDF file) OR
%Error 1000 An error occurred while parsing a contents stream. Unable to analyze the PDF file.
%This is a known problem with pdfLaTeX conversion filter. The file cannot be opened with acrobat reader
%Please use one of the alternatives below to circumvent this error by uncommenting one or the other
%\pdfobjcompresslevel=0
%\pdfminorversion=4

% See the \addtolength command later in the file to balance the column lengths
% on the last page of the document

% The following packages can be found on http:\\www.ctan.org
\usepackage{graphics} % for pdf, bitmapped graphics files
\usepackage{epsfig} % for postscript graphics files
\usepackage{mathptmx} % assumes new font selection scheme installed
\usepackage{times} % assumes new font selection scheme installed
\usepackage{amsmath} % assumes amsmath package installed
\usepackage{amssymb}  % assumes amsmath package installed
\usepackage{subcaption}

\title{\LARGE \bf

Single Source Synergy-Constrained Policy Learning for Multi-finger In-Hand Manipulation Across Novel Objects with Tactile Sense
}


%\author{Serena Trovalusci$^{1}$ and Mitsuhiro Hayashibe$^{2}$% <-this % stops a space
%\thanks{*This work was supported by Sapienza University of Rome and Tohoku University}% <-this % stops a space
%\thanks{$^{1}$S. Trovalusci was with the Department of Computer, Control, and 
        %Management Engineering, Sapienza University of Rome, 00185 Rome, Italy.
        %{\tt\small trovalusciserena@gmail.com}}%
%\thanks{$^{2}$M. Hayashibe is with the Department of Robotics, Graduate School 
        %of Engineering, Tohoku University, Sendai 980-8579, Japan.
        %{\tt\small mitsuhiro.hayashibe.e6@tohoku.ac.jp}}%



\begin{document}



\maketitle
\thispagestyle{empty}
\pagestyle{empty}


%%%%%%%%%%%%%%%%%%%%%%%%%%%%%%%%%%%%%%%%%%%%%%%%%%%%%%%%%%%%%%%%%%%%%%%%%%%%%%%%
\begin{abstract}
        % Dexterous in-hand manipulation with anthropomorphic robotic hands remains 
        % an open challenge in robotics, largely due to the sample inefficiency of 
        % reinforcement learning in high-dimensional continuous action spaces. 
        % Inspired by the motor synergy hypothesis in neuroscience, we propose the 
        % Motor Synergy Generalization Framework, a three-stage pipeline in which 
        % spatial synergies are extracted offline via Principal Component Analysis 
        % from the joint-action trajectories of a converged full-DoA policy on a 
        % source task, and reused as a fixed low-dimensional action decoder for 
        % training new agents on novel target tasks. We evaluate the framework on 
        % five object reorientation tasks with substantially different geometries 
        % on the Shadow Dexterous Hand in MuJoCo, comparing a synergy-constrained 
        % agent operating in a 5-dimensional latent space against a full-action-space 
        % baseline in the native 20-dimensional joint space. The synergy-constrained 
        % agent consistently reaches comparable or higher success rates in fewer 
        % environment timesteps and less wall-clock training time across all five 
        % objects, with speedups ranging from $1.5\times$ on the large cube to 
        % $6\times$ on the sphere, demonstrating that a synergy basis extracted from 
        % a single source task acts as a geometry-agnostic motor prior for dexterous 
        % manipulation.
Multi-finger dexterous in-hand manipulation remains a grand challenge in robotics due to high-dimensional action spaces and complex contact dynamics. While Deep Reinforcement Learning (DRL) offers a viable data-driven solution, its practical application is severely bottlenecked by extreme sample inefficiency during random exploration. Inspired by biological motor synergies, this paper proposes a novel three-stage pipeline that extracts spatial coordination patterns from a single source manipulation task to accelerate learning across a wide range of novel object geometries. First, a full-Degree-of-Actuation (DoA) Soft Actor-Critic (SAC) agent combined with Hindsight Experience Replay (HER) is trained on a standard cube reorientation task. Second, spatial synergies are extracted offline via Principal Component Analysis (PCA) from the converged policy's joint-action trajectories. Third, the resulting 5-dimensional synergy basis is frozen and reused identically as a low-dimensional action decoder to constrain policy learning on novel target tasks. We evaluate our approach on five target objects with substantially different shapes: a small cube, a large cube, an egg, a cylinder, and a sphere. Simulation experiments on the 24-Degree-of-Freedom (DoF) Shadow Hand demonstrate that the synergy-constrained agent outperforms the full-DoA baseline, reducing sampling duration and wall-clock training time by up to 67\% on the spherical geometry. Our findings demonstrate that a low-dimensional kinematic synergy basis extracted from a single source object can act as a geometry-agnostic motor prior capable of accelerating dexterous manipulation learning for unseen geometries, in addition this advantage is robust to the lack of tactile sensing.
\end{abstract}


%%%%%%%%%%%%%%%%%%%%%%%%%%%%%%%%%%%%%%%%%%%%%%%%%%%%%%%%%%%%%%%%%%%%%%%%%%%%%%%%
\section{INTRODUCTION}

\begin{figure*}[h]
        \centering
        \begin{subfigure}[b]{0.18\textwidth}
            \includegraphics[width=\textwidth]{figures/1.png}
            \caption{}
        \end{subfigure}
        \hfill
        \begin{subfigure}[b]{0.18\textwidth}
            \includegraphics[width=\textwidth]{figures/2.png}
            \caption{}
        \end{subfigure}
        \hfill
        \begin{subfigure}[b]{0.18\textwidth}
            \includegraphics[width=\textwidth]{figures/3.png}
            \caption{}
        \end{subfigure}
        \hfill
        \begin{subfigure}[b]{0.18\textwidth}
            \includegraphics[width=\textwidth]{figures/4.png}
            \caption{}
        \end{subfigure}
        \hfill
        \begin{subfigure}[b]{0.18\textwidth}
            \includegraphics[width=\textwidth]{figures/5.png}
            \caption{}
        \end{subfigure}
        \caption{Time-lapse sequence of the source task: goal-conditioned in-hand reorientation of a standard cube using the 20-DoA hand. The policy dynamically manipulates the object from its initial configuration (a), through coordinated finger-gaiting transitions (b–d), to align with the target configuration (e). 
        Tactile readings are distributed across the hand surface. }
        \label{fig:cube_sequence}
\end{figure*}

In-hand manipulation is the ability to reorient and reposition an object within 
the hand without setting it down or using the other hand. It is one of the most 
fundamental yet complex motor capabilities in humans, as it requires the 
coordinated motion of fingers and contact forces, with sensory feedback. Current industrial and medical robotic hands already achieve remarkable precision in structured and repetitive tasks. The next frontier in robotics is dexterous in-hand manipulation in unstructured environments: the ability to manipulate arbitrary objects with the kind of fluid adaptability that humans exercise effortlessly every day. Performing this with anthropomorphic robotic hands remains one of the most open and challenging problems in robotics \cite{c1, c20}.

The intricate architecture that characterizes the human hand with its bones, muscles, and nerves makes in-hand manipulation tasks highly redundant. Previous studies have revealed that the Central Nervous System (CNS) does not control every single Degree of Freedom (DoF) independently; instead, movements are generated by a small number of functional modules, called \textit{motor synergies} \cite{c2}. A relevant feature of these synergies is their reusability across tasks, which allows the CNS to maintain a dictionary of reusable coordination primitives and learn new behaviours by re-using and combining them \cite{c3,c4}. 

In robotics, rigid anthropomorphic hands such as the Shadow Dexterous Hand involve up to 
24 DoF, of which 20 are independently actuated. This makes 
model-based control extremely challenging due to 
high-dimensional action spaces and non-linear contact dynamics. Data-driven approaches, such as Deep Reinforcement Learning (DRL), have shown promising results in recent years. Rather 
than requiring an explicit model of the environment or an expensive dataset of 
expert demonstrations \cite{c23}, DRL allows the agent to learn optimal policies directly 
from trial-and-error interactions with the environment, guided only by a reward 
signal \cite{c6}. However, random exploration in high-dimensional continuous 
action spaces is extremely slow, and training requires enormous numbers of 
environment interactions. This sample inefficiency remains the central bottleneck 
limiting the practical applicability of DRL to dexterous in-hand manipulation \cite{c7}. Computationally, seeking motor synergy combinations can be seen as a 
dimensionality reduction technique, which reduces the complexity 
of the control problem and enables the rapid acquisition of complex motor 
skills \cite{c5}, provided the resulting low-dimensional controller is 
designed to exploit the natural dynamics of the system being 
controlled \cite{c22}.

Inspired by the Synergistic Action Representation (SAR) of Berg et al.~\cite{c9}, 
which showed that a muscle synergy basis distilled from multiple manipulation 
tasks transfers zero-shot to novel objects, and also by
the motor synergy generalization framework~\cite{c8}~\cite{c19}, which demonstrated that arm 
synergies extracted from DRL policies trained on multiple reaching targets can 
be combined to generalize to entirely new targets, this paper investigates whether 
motor synergies extracted from a single source in-hand manipulation task can 
accelerate reinforcement learning on a set of novel target tasks 
involving objects with substantially different geometries.
We propose a three-stage pipeline in which a full-Degree-of-Actuation (DoA) 
Soft Actor-Critic (SAC) agent is first trained on a source task (Stage~A), spatial 
synergies are then extracted offline via Principal Component Analysis (PCA) from 
the joint-action trajectories of the converged policy (Stage~B), and the resulting 
synergy basis is finally used as a fixed low-dimensional action decoder for 
training new DRL agents on target tasks (Stage~C). The specific contributions of 
this work are as follows:
\begin{itemize}
    \item We propose a three-stage 
    pipeline that extracts spatial synergies from a single source policy and 
    reuses them as a fixed action decoder for novel target tasks on the Shadow 
    Dexterous Hand.
    \item We demonstrate that constraining the action space to a 5-dimensional 
    synergy subspace extracted from cube manipulation accelerates learning on five 
    target tasks spanning a wide range of object geometries, in terms of both 
    sample efficiency and wall-clock training time.
    \item We show that the synergy advantage is consistent across all evaluated 
    objects despite the synergy basis having been extracted from a single source 
    task with a qualitatively different geometry from most of the targets.
\end{itemize}

\section{RELATED WORK}

The neuroscientific observation that the human hand operates through a 
low-dimensional set of coordinated joint patterns was first formalised in 
a robotic context by Santello, Flanders and Soechting~\cite{c3}, who showed 
that the first two principal components of hand posture data account for 
more than 80\% of the variance across a wide variety of grasping tasks, 
suggesting that a small number of coordination patterns underlies the full 
repertoire of human grasping. Todorov and Ghahramani~\cite{c21} similarly analyzed postural synergies extracted via PCA during complex manipulation tasks, extending this line of work beyond simple grasping. Ciocarlie and Allen~\cite{c10} 
exploited this result for autonomous grasp planning, reducing an otherwise 
high-dimensional search in configuration space to a two-dimensional synergy 
plane with negligible loss in grasp quality. Bicchi and collaborators 
extended the idea to mechanical design through the theory of soft synergies 
and synergy-based underactuation~\cite{c11, c12}, yielding the Pisa/IIT SoftHand, 
a single-motor hand capable of grasping a wide variety of objects by virtue 
of its mechanically realised first synergy.

Chai and Hayashibe~\cite{c13} were the first to show that high-performing 
DRL policies spontaneously develop motor synergy structures, and that 
the degree of synergy structure correlates with task performance. It is extended to full-dimensional arm manipulation~\cite{c14}, 
confirming that synergy emergence is a robust phenomenon in high-dimensional 
continuous action spaces and that it can be exploited to improve learning 
efficiency. The work directly related to the present paper is that of Kutsuzawa and Hayashibe~\cite{c8}, who extracted spatio-temporal synergies from Proximal Policy Optimization (PPO) policies
via Non-negative Matrix Factorization (NMF) and demonstrated that combining synergy repertoires from two 
single-plane reaching policies generalizes to entirely new planes that 
neither repertoire could reach alone. Our work inherits this core insight 
and extends it to the harder setting of multi-finger in-hand manipulation 
with sparse rewards, replacing NMF with PCA and PPO with SAC combined with 
Hindsight Experience Replay (HER). At the muscle level, Berg, Caggiano and Kumar~\cite{c9} developed the
SAR, obtained by composing Independent Component Analysis (ICA) and PCA 
on muscle activations produced by a play-phase policy on simple tasks. The 
resulting basis generalizes zero-shot to in-hand reorientation of a hundred 
novel objects, without any additional training on the target tasks. In a 
smaller-scale, single-source experiment, Berg et al. also extract a 
\emph{kinematic} synergy basis directly on the 20-DoA Shadow Dexterous Hand 
used in this work, from a policy trained on a single object, an egg, and 
use it to accelerate training on two 
unseen geometries: a capsule and a cylinder. Our work extends this 
single-source, kinematic setting considerably further: rather than 
transferring within a family of smooth, edgeless shapes, we extract 
synergies from a single cube and evaluate generalization across 
five geometries spanning a substantially wider range, including a sphere, 
its near-opposite in contact structure. We do so within a full SAC+HER 
training loop under sparse binary rewards, which allows the 
policy to adapt the temporal activation of the frozen synergy basis to 
each target object geometry.
\begin{figure*}[t]
        \centering
        \includegraphics[width=\textwidth]{figures/method_overview_1.png}
        \caption[]{Framework overview: Stage~A trains a full-DoA
        SAC agent on the source task. Stage~B extracts the synergy basis
        offline via PCA. Stage~C trains a synergy-constrained SAC agent
        on each target task using the frozen decoder.}
        \label{fig:method-overview}
\end{figure*}

Andrychowicz et al.~\cite{c7} demonstrated that DRL policies trained entirely 
in simulation on the Shadow Dexterous Hand can transfer to a real robot 
using domain randomisation, operating in the full 20-dimensional joint 
space at the cost of enormous computational resources. Plappert et 
al.~\cite{c16} introduced the multi-goal manipulation benchmarks on the Shadow 
Hand used throughout this paper, providing the first systematic evaluation 
of Deep Deterministic Policy Gradient (DDPG) combined with HER on in-hand reorientation tasks. HER~\cite{c17} addresses the fundamental difficulty of learning 
from sparse binary rewards in high-dimensional spaces by relabeling failed 
trajectories with goals that were actually achieved, recovering a dense 
learning signal from otherwise uninformative episodes. For policy 
optimisation, we adopt SAC~\cite{c15}, an off-policy 
maximum-entropy algorithm that has become the standard for 
continuous-control benchmarks due to its sample efficiency and principled 
exploration.

The present work occupies a specific and previously unexplored position in 
this space. Like Santello et al. and Ciocarlie and Allen, we study 
postural-level kinematic synergies extracted via PCA. We learn these synergies from DRL agent trajectories rather 
than human data, and we freeze the basis for transfer to novel tasks. 
Unlike any of the above, we combine these ingredients in the specific 
setting of multi-finger in-hand reorientation on the Shadow Dexterous Hand 
with sparse binary rewards and HER, demonstrating that a kinematic synergy 
basis extracted from a single source object accelerates learning on objects 
whose geometry was never seen during synergy extraction.


\section{METHODS}
\label{sec: methods}

We address the task of \textit{dexterous in-hand reorientation}: an object is placed on the palm of the Shadow Dexterous Hand. It must grasp and rotate the object from a randomised initial orientation to a randomised target orientation, without releasing it, as 
shown in Fig.~\ref{fig:cube_sequence}. The task is formulated as a 
multi-goal Markov Decision Process \cite{c16}: the policy $\pi(a|s,g)$ is conditioned on both the state $s$ and the goal $g$, and the reward is computed through a goal-parametrised function $r(s,g)$ rather than a single fixed reward. This decoupling of state and goal is what allows HER to relabel a failed trajectory with a goal it actually achieved, recovering a learning signal from otherwise uninformative episodes. At the beginning of each episode a
target orientation is sampled uniformly from SO(3), and the agent must
reach it within a rotational tolerance of $\epsilon_r = 0.1$\,rad and a
positional tolerance of $\epsilon_p = 1$\,cm. The reward is sparse and
binary: the agent receives $0$ upon success and $-1$ otherwise. The
control frequency is 25\,Hz, and each episode lasts a fixed $T = 100$
simulation steps regardless of task completion, as these tasks are non-time-dependent continuing tasks with no early termination on success. Following the hand-manipulation environments of \cite{c16}, the observation includes the 24 joint positions and velocities of the Shadow Dexterous Hand, together with the manipulated object's Cartesian position and quaternion orientation, and its linear and angular velocities. We consider two observation conditions: a base condition using this state alone, and a touch-sensor condition that augments it with 92 continuous tactile readings distributed across the hand surface.

We report two complementary success criteria throughout this paper.
Under the \emph{sustained-hold} criterion, an episode is successful if
the achieved and target orientations are within tolerance at the
episode's final timestep; this is the primary criterion used for the
learning curves and sample/wall-clock efficiency comparisons in
Section~\ref{sec:experiments}. Under the \emph{reorientation} criterion, an episode is
successful if the tolerance is satisfied at any timestep during the
episode, regardless of whether it is subsequently maintained; we report
this as a complementary metric to characterise how often the target
orientation is reached at all, independently of the agent's ability to
stabilise the object once there.

We consider six object geometries in total: one for the \textit{source task} and five for the \textit{target tasks}. The \textit{source task} uses a cube of side $0.025$\,m. The five \textit{target tasks} include two scaled variants of the source cube, a small cube ($0.020$\,m side) and a large cube ($0.030$\,m side), and three shapes with substantially different contact geometry from the cube: an egg (ellipsoid, $r_x = r_y = 0.030$\,m, $r_z = 0.040$\,m), a cylinder (radius $0.030$\,m, height $0.050$\,m), and a sphere (radius $0.033$\,m). The small and large cubes were obtained by modifying the object geometry in the MuJoCo XML of the source cube, the egg environment is used unmodified from Gymnasium-Robotics \cite{c18}, and the cylinder and sphere environments were created for this work.

We propose a three-stage pipeline, illustrated in 
Fig.~\ref{fig:method-overview}, whose key design principle is that the 
synergy basis extracted in Stage~B is never re-fitted or fine-tuned on 
the target tasks: the same low-dimensional motor vocabulary is reused 
identically across all target objects. Both the full-action-space 
baseline and the synergy-constrained agent share the same network 
architecture: actor and critics are multi-layer perceptrons with two 
hidden layers of 256 units and ReLU activations. The only architectural 
difference is the output dimension of the actor (20 for the baseline, 
5 for the synergy-constrained agent) and the corresponding input 
dimension of the critics. All other hyperparameters are kept identical, 
so that any observed difference in learning performance can be attributed 
solely to the action-space parameterisation.

\subsection{Stage A: Full-DoA source task training}

In Stage A, the agent is trained on the source task acting directly in the 20-dimensional joint command space of the Shadow Dexterous Hand. The cubic geometry of the source object, with flat faces, sharp edges, 
and rotational symmetry, requires the hand to produce a rich and 
varied repertoire of multi-finger coordination patterns across arbitrary 
target orientations. The resulting joint-action trajectories consequently span a wide range of coordinated finger movements that can be reused across all target tasks. The Stage~A agent is trained for $6 \times 10^6$ environment
timesteps, reaching a sustained-hold success rate of approximately 50\%
and a reorientation success rate of 95\% over 100 evaluation episodes. The sustained-hold success rate reflects the policy's ability to stabilise the object at the target orientation for the remainder of the episode, whereas the reorientation success rate reflects its ability to reach the target orientation at all. For the purpose of synergy extraction, it is the latter that matters, since the synergy basis is distilled from the joint-action trajectories generated while approaching the target orientation rather than from the agent's ability to sustain the final pose. We therefore consider this policy sufficient for synergy extraction despite its moderate sustained-hold success rate.


\subsection{Stage B: Synergy extraction}

Once the Stage~A policy training is complete, it is frozen and executed for 
$N = 200$ evaluation rollouts of length $T = 100$ timesteps each. At 
every timestep we record the commanded joint action 
$a_t \in \mathbb{R}^{20}$, yielding a data matrix 
$X \in \mathbb{R}^{20{,}000 \times 20}$.

Spatial synergies are extracted from $X$ via PCA. The centred matrix 
$\tilde{X} = X - \mathbf{1}\bar{x}^\top$ is decomposed via SVD, and 
the top-$K$ right singular vectors form the rows of the synergy matrix 
$W \in \mathbb{R}^{K \times 20}$. Each row $w_k$ is a spatial synergy: 
a coordinated pattern of joint displacements that can be linearly 
combined with others to approximate the original joint commands. The 
reconstruction of a single action sample reads
%
\begin{equation}
    a_t \approx \bar{x} + \sum_{k=1}^{K} h_{tk}\, w_k = W^\top h_t + \bar{x},
    \label{eq:decoder}
\end{equation}
%
where $h_t \in \mathbb{R}^K$ are the synergy activations and $\bar{x}$ 
is the mean action vector.

The number of retained components $K$ is chosen by combining reconstruction quality with task-performance evaluation. We compute the mean per-DoA coefficient of determination $\overline{R^2}(K)$ on a held-out test set, splitting 20\% of the rollouts by trajectory to avoid temporal leakage. Reconstruction quality rises steeply over the first few components and then saturates slowly: $17$ components are needed to reach around $\overline{R^2} = 0.95 $, and all the $20$ components are required for exact reconstruction. Saturating reconstruction is not the right objective, since retaining 17 out of $20$ dimensions would defeat the purpose of the method. We therefore retain only $5$ components, which reach around $\overline{R^2} = 0.5 $ and prove entirely sufficient to learn all target tasks with a substantial advantage over the baseline.


\subsection{Stage C: Synergy-constrained target task training}

In Stage~C, a new SAC+HER agent is trained on each target task. The 
policy $\pi_\theta(h \mid s, g)$ outputs a $K$-dimensional synergy 
activation vector $h_t \in \mathbb{R}^K$ instead of a direct joint 
command. At every control step, the activation is decoded into a full 
20-dimensional joint command via the fixed affine map of 
\eqref{eq:decoder}, clipped to the simulator's native bounds $[-1, 1]$ 
and applied to the hand. The decoder parameters $W$ and $\bar{x}$ 
receive no gradient updates during Stage~C; they are loaded once from 
Stage~B and remain frozen for every target task.

The synergy activation space is bounded by $[-h_{\max}, h_{\max}]$ with 
$h_{\max} = 3.0$. Setting $h_{\max} = 1.0$ to match the native action 
bounds proved insufficient, as the PCA activations span a range 
considerably larger than $[-1, 1]$; expanding to $h_{\max} = 3.0$ 
resolved poor learning progress across all objects.

\section{EXPERIMENTS}
\label{sec:experiments}
\subsection{Experimental setup}

\begin{table}[b]
        \caption{SAC+HER hyperparameters}
        \label{tab:hyperparam}
        \begin{center}
        \begin{tabular}{ll}
        \hline
        \textbf{Hyperparameter} & \textbf{Value} \\
        \hline
        Learning rate (actor and critics) & $3 \times 10^{-4}$ \\
        Batch size & 256 \\
        Discount factor $\gamma$ & 0.99 \\
        Replay buffer size & $10^6$ \\
        HER strategy & future \\
        HER relabelled goals & 8 \\
        Entropy target & $-0.8 \times \dim(\mathcal{A})$ \\
        Synergy activation range $h_{\max}$ & 3.0 \\
        Hidden layers & {[256, 256]} \\
        Training timesteps & $3 \times 10^6$ \\
        Parallel environments & 16 \\
        Gradient steps per update & 4 \\
        Training device & CPU (24 cores) \\
        \hline
        \end{tabular}
        \end{center}
        \end{table}


All experiments are conducted in simulation using MuJoCo, on top of the in-hand manipulation environments from Gymnasium-Robotics~\cite{c18}. The synergy basis is extracted once from the cube source task and reused without modification across all five targets. For each target task, we train and evaluate under both observation conditions introduced in Section \ref{sec: methods}: the touch-sensor condition, which augments the observation with 92 continuous tactile readings distributed across the hand surface, and a base condition in which these readings are removed. 

For each target task and observation condition, two agents are trained: the \emph{full-action-space} baseline, operating in the native 20-dimensional joint space, and the \emph{synergy-constrained} agent, operating in the 5-dimensional synergy activation space. All hyperparameters are identical between the two agents and are listed in Table~\ref{tab:hyperparam}. Each agent is trained for $3 \times 10^6$ environment timesteps with SAC combined with HER, and every configuration is repeated with 3 independent random seeds.

Success rate is estimated under the two definitions introduced in Section~\ref{sec: methods}. The sustained-hold success rate is logged automatically during training via deterministic evaluation over 50 episodes every $10^5$ timesteps, and is the metric used for the sample and wall-clock efficiency comparisons below. The reorientation success rate is estimated post-hoc, by re-evaluating the periodic training checkpoints (saved every $5\times10^5$ timesteps) over 50 deterministic episodes each. We report it as a complementary, descriptive metric characterising how often the target orientation is reached at all, independently of the agent's ability to subsequently stabilise the object.


\subsection{Results}

\begin{figure*}[t]
    \centering
    \includegraphics[width=\textwidth]{figures/panels/large_cube.pdf}\\[-4pt]
    \includegraphics[width=\textwidth]{figures/panels/egg.pdf}\\[-4pt]
    \includegraphics[width=\textwidth]{figures/panels/cylinder.pdf}\\[-4pt]
    \includegraphics[width=\textwidth]{figures/panels/sphere.pdf}
    \caption{Sustained-hold and reorientation success rate (mean $\pm$ std over 3 seeds), with and without touch sensors, for each evaluated target object.}
    \label{fig:curves-all}
\end{figure*}


\begin{table}[t]
    \caption{Sampling duration to sustained-hold threshold (timesteps, $\times 10^6$). BL = baseline, SC = synergy-constrained. Values are mean $\pm$ std across seeds that cross the threshold; "--" = not reached by all seeds within budget.}
    \label{tab:steps}
    \begin{center}
    \begin{tabular}{llcc}
    \hline
    \textbf{Object} & \textbf{Sensors} & \textbf{BL} & \textbf{SC} \\
    \hline
    Large cube & yes & -- & $2.03 \pm 0.24$ \\
    Large cube & no  & -- & $2.63 \pm 0.33$ \\
    Egg        & yes & $2.23 \pm 0.17$ & $1.57 \pm 0.13$ \\
    Egg        & no  & -- & $1.80 \pm 0.22$ \\
    Cylinder   & yes & -- & $1.87 \pm 0.13$ \\
    Cylinder   & no  & -- & $1.67 \pm 0.13$ \\
    Sphere     & yes & $1.87 \pm 0.13$ & $0.77 \pm 0.13$ \\
    Sphere     & no  & $2.63 \pm 0.25$ & $0.87 \pm 0.17$ \\
    \hline
    \end{tabular}
    \end{center}
\end{table}


\begin{table}[t]
    \caption{Wall-clock time to sustained-hold threshold. BL = baseline, SC = synergy-constrained. Values are mean $\pm$ std across seeds that cross the threshold; "--" = not reached by all seeds within budget.}
    \label{tab:time}
    \begin{center}
    \begin{tabular}{llcc}
    \hline
    \textbf{Object} & \textbf{Sensors} & \textbf{BL} & \textbf{SC} \\
    \hline
    Large cube & yes & -- & $1$h$13$m $\pm 8$m \\
    Large cube & no  & -- & $1$h$27$m $\pm 10$m \\
    Egg        & yes & $1$h$19$m $\pm 6$m & $0$h$57$m $\pm 4$m \\
    Egg        & no  & -- & $1$h$00$m $\pm 6$m \\
    Cylinder   & yes & -- & $1$h$07$m $\pm 4$m \\
    Cylinder   & no  & -- & $1$h$05$m $\pm 7$m \\
    Sphere     & yes & $1$h$08$m $\pm 5$m & $0$h$29$m $\pm 4$m \\
    Sphere     & no  & $1$h$30$m $\pm 9$m & $0$h$30$m $\pm 5$m \\
    \hline
    \end{tabular}
    \end{center}
\end{table}


Fig.~\ref{fig:curves-all} reports the sustained-hold and reorientation success-rate curves for each evaluated target task and observation condition (mean $\pm$ std over 3 seeds). Sample efficiency and wall-clock training duration are computed from the sustained-hold curves using a threshold set to 80\% of the peak mean success rate attained by either agent. For each seed individually, we record the training checkpoint at which its own sustained-hold curve first crosses this threshold, and report the mean and standard deviation of these per-seed crossing points only when all three seeds cross within the training budget. In several object/condition combinations, one or more seeds of the full-action-space baseline never cross the threshold at all; we report these cases explicitly as "not reached".

The large cube proves the most challenging target overall: even the synergy-constrained agent's sustained-hold success rate peaks at only 18.7\% (touch-sensor) and 16.7\% (no-sensor). The qualitative gap nonetheless persists: as Tables~\ref{tab:steps} and~\ref{tab:time} show, the synergy-constrained agent reaches even this low threshold under both observation conditions, while every seed of the baseline fails to do so within the full training budget. The reorientation curves in Fig.~\ref{fig:curves-all} offer a more encouraging picture, however: the synergy-constrained agent's reorientation success rate reaches approximately 70\% (touch-sensor) and 62\% (no-sensor), indicating that the large cube exposes a genuine difficulty in stabilising the object once reoriented, rather than an inability to reorient it in the first place.

On the egg, both agents reach a high sustained-hold peak (98.0\% touch-sensor, 94.7\% no-sensor), but the synergy-constrained agent reaches the corresponding threshold considerably faster (Tables~\ref{tab:steps},~\ref{tab:time}) -- approximately 30\% fewer samples and 28\% less wall-clock time under the touch-sensor condition. Removing tactile observation widens this gap further: the synergy-constrained agent still reaches its threshold, while every seed of the baseline fails to do so within the full training budget.

On the cylinder, the synergy-constrained agent reaches its sustained-hold threshold under both observation conditions (Tables~\ref{tab:steps},~\ref{tab:time}), while every seed of the baseline fails to do so within the full training budget in either condition. The cylinder is thus the clearest instance in our results of a capability gap within a fixed compute budget: the baseline does not reach a level of sustained performance that the synergy-constrained agent attains in roughly half the allotted training.

The sphere yields the cleanest and most striking result. Both agents saturate to 100\% sustained-hold success within the training budget under both observation conditions, so the standard absolute 80\% threshold applies throughout. The synergy-constrained agent reaches it roughly $2.4\times$ faster than the baseline with tactile sensing, and roughly $3\times$ faster without it (Tables~\ref{tab:steps},~\ref{tab:time}) -- reductions of up to 67\% in both sampling duration and wall-clock time. We attribute the sphere's outsized speedup, and its robustness to the loss of tactile information for the synergy-constrained agent specifically, to the high symmetry of the spherical geometry: an identical contact surface is presented in every orientation, making the reorientation task intrinsically simpler and the synergy basis an immediately effective exploration prior even without touch feedback.

Across the five target objects evaluated so far, the synergy-constrained agent consistently reaches a given sustained-hold performance level in fewer environment interactions and less wall-clock training time than the full-action-space baseline, and in several object/observation combinations the baseline fails to reach that level at all within the fixed training budget. This advantage holds under both the touch-sensor and no-sensor conditions, and in most cases widens when tactile information is removed, suggesting that the synergy basis provides a structural prior that partially compensates for impoverished sensory feedback.



\section{DISCUSSION}
\subsection{Mechanisms of Cross-Geometry Generalization}
The central and most notable finding of this work is that the synergy basis extracted from the cubic geometry transfers effectively to objects with entirely different topologies and contact profiles, such as the egg and the cylinder. We attribute this robust generalization to the intrinsic properties of the source task. The standard cube possesses sharp edges, flat faces, and vertices, requiring the 20-DoA (Degree of Actuation) hand to generate a highly diverse dictionary of coordinated multi-finger behaviors. 

Consequently, the PCA in Stage B does not merely memorize "cube-specific" poses. Instead, it distills the fundamental kinematic primitives of anthropomorphic manipulation, such as coordinated finger flexion/extension and opposing thumb movements. When transferred to an egg or a cylinder, these primitives serve as a highly effective exploration prior. While the baseline agent in the full 20-dimensional action space wastes millions of samples exploring destructive, uncoordinated joint movements that often lead to dropping the object, the synergy-constrained agent is restricted to exploring within a subspace of naturally coordinated movements, drastically reducing the search space while maintaining the capacity to manipulate novel geometries.


\subsection{Trade-off Between Dimensionality and Policy Expressivity}
In Stage B, we selected $K=5$ principal components, which captured $\overline{R^2}=0.512$ of the reconstruction variance in the source trajectory data. From a purely kinematic reconstruction standpoint, this is relatively low: seventeen components are needed to reach $\overline{R^2}=0.946$, and the full twenty are required for exact reconstruction. However, our results demonstrate a key paradox in synergy-based control: maximizing spatial reconstruction accuracy is not equivalent to maximizing reinforcement learning efficiency. 

Retaining 17 dimensions does not constitute a substantial reduction in dimensionality and would defeat the purpose of the method. Conversely, restricting the action space to $K=5$ acts as a powerful regularizer. Although a fixed 5-dimensional spatial decoder introduces geometric reconstruction errors, we hypothesize that the online DRL agent in Stage C overcomes this spatial limitation through \textit{temporal adaptation}. Because the SAC policy can dynamically modulate the activations $h_t$ at 25 Hz, it learns to sequentially combine the frozen spatial primitives over time to fit the unique curvatures of the target objects. This strongly suggests that the time-varying combination of a small, rigid spatial vocabulary is sufficient for highly adaptive dexterous manipulation.

\subsection{Robustness to Impoverished Sensory Feedback}
Beyond geometry generalisation, our ablation over the touch-sensor observation reveals a further property of the synergy prior: when the 92-dimensional tactile observation is removed, the advantage over the full-action-space baseline increases. Across the five objects evaluated under this ablation, the baseline agent's sample efficiency degrades substantially more than the synergy-constrained agent's when tactile feedback is unavailable, and in several cases the baseline fails to reach a usable level of performance at all within the training budget without it. We interpret this as evidence that the synergy basis itself encodes part of the exploratory function normally provided by rich sensory feedback. The frozen decoder already restricts exploration to naturally coordinated finger movements, and the synergy-constrained agent depends less on tactile information to discover effective contact strategies, whereas the full-action-space baseline must discover both the coordination pattern and the appropriate contact strategy from tactile feedback simultaneously.
This is consistent with Berg et al.~\cite{c9}, who report that sight and touch ``were not required to yield successful control policies'' for their synergy-based models and who separately obtained their own Shadow Hand results without tactile sensing.


\subsection{Limitations and Boundary Conditions}
Despite the robust generalization achieved across the evaluated shapes, several boundary conditions must be acknowledged. First, the spatial decoder extracted via PCA remains entirely linear and static throughout Stage C. While temporal modulation mitigates spatial inaccuracies, a rigid linear subspace inherently bounds the ultimate expressivity of the hand. For highly complex, asymmetric, or non-convex objects, a fixed 5-dimensional linear prior derived from a symmetric cube may over-regularize the system, preventing the generation of highly specific behaviors needed to manipulate tools with distinct handles or fragile articulated mechanisms. 

Second, task difficulty varies substantially across the evaluated geometries: on the large cube, the synergy-constrained agent's sustained-hold success rate does not exceed approximately 16--19\% within the $3\times10^6$-timestep training budget (depending on the observation condition), even though its reorientation success rate reaches approximately 60--70\%, indicating that the object can be reliably reoriented but not consistently stabilised there within the given budget. We suspect this reflects the specific geometric relationship between the large cube and the source cube used for synergy extraction -- although the two share the same rotationally symmetric, flat-faced topology, the increased size may shift the effective range of finger-to-surface contact distances outside what the source policy's synergy repertoire was tuned for -- but confirming this hypothesis, and closing this stabilisation gap, is left to future work.

Third, our approach still requires a full training loop on each target task; the synergy basis accelerates learning but does not eliminate the need for target-specific training.

Fourth, our current pipeline is evaluated within a simulated environment (MuJoCo). While the use of 92 continuous tactile readings provides rich sensory feedback, simulation models often simplify critical real-world phenomena such as soft-pad deformation, material friction variance, and actuator backlash. Transitioning the synergy-constrained policy to a physical platform without further domain randomization or online calibration could lead to performance degradation due to the sim-to-real gap.

\subsection{Future Work}
To address these limitations, our immediate future work will focus on introducing \textit{online synergy adaptation}. Instead of completely freezing the parameters $W$ and $\overline{x}$ in Stage C, we plan to implement a slow-moving meta-learning loop or parameter-efficient fine-tuning (PEFT) layers to allow the low-dimensional motor vocabulary to slowly warp and fit highly anomalous target geometries. 

Additionally, we aim to transition from a purely linear PCA mapping to non-linear action decoders, such as deep variational autoencoders (VAEs), to capture more complex, non-linear joint correlations while keeping the action space low-dimensional. Finally, we intend to deploy the learned policies on a physical robotic hand. By incorporating robust domain randomization over tactile noise and contact friction during the Stage C training, we expect the synergy-constrained policy to exhibit superior resilience during physical sim-to-real transfer compared to full-DoA alternatives. 




%%%%%%%%%%%%%%%%%%%%%%%%%%%%%%%%%%%%%%%%%%%%%%%%%%%%%%%%%%%%%%%%%%%%%%%%%%%%%%%%



%%%%%%%%%%%%%%%%%%%%%%%%%%%%%%%%%%%%%%%%%%%%%%%%%%%%%%%%%%%%%%%%%%%%%%%%%%%%%%%%
\section{CONCLUSIONS}
In this work, we presented a novel three-stage pipeline for single-source synergy-constrained policy learning to address the critical sample inefficiency bottleneck in dexterous multi-finger in-hand manipulation. By training a full-DoA Soft Actor-Critic agent on a single standard cube reorientation task, extracting spatial synergies offline via Principal Component Analysis, and freezing the resulting $K=5$ dimensional synergy basis as a static action decoder, we successfully constrained the exploration space for novel target objects without any target-specific fine-tuning of the decoder parameters.  Our empirical evaluations on the 24-DoF Shadow Dexterous Hand across five geometrically diverse target tasks -- encompassing scaled cubes, an ellipsoid egg, a cylinder, and a sphere, under both tactile and tactile-free observation conditions -- demonstrate that this low-dimensional kinematic prior consistently and significantly accelerates DRL. The proposed approach yielded substantial reductions in both sampling duration and wall-clock training time across all novel shapes, reducing both by up to 67\% on the spherical geometry. This advantage was found to be robust to the removal of tactile sensing and, in most cases, amplified by it.   Crucially, these findings demonstrate that coordination primitives distilled from a single source object can generalize effectively as geometry-agnostic motor priors. This validates the bio-inspired hypothesis that a compact, rigid spatial vocabulary, when paired with high-frequency temporal modulation, provides sufficient expressivity to manipulate unseen physical structures. As single source hand synergy can be practically obtained, this approach also contributes to reduce burden for large real data collection for robot hand learning.
Ultimately, this framework offers a highly efficient and scalable paradigm for training high-dimensional anthropomorphic hands, bringing us a step closer to achieving fluid, human-like dexterity in unstructured environments. 


\begin{thebibliography}{99}

        \bibitem{c1}
        Y. Huang, D. Fan, H. Duan, D. Yan, W. Qi, J. Sun, Q. Liu, and P. Wang, ``Human-like dexterous manipulation for anthropomorphic five-fingered hands: A review,'' \emph{Biomimetic Intelligence and Robotics}, vol. 5, no. 1, Art. no. 100212, 2025.

        \bibitem{c2}
        N. A. Bernstein, ``The coordination and regulation of movements,'' \emph{Pergamon Press}, 1967.
        
        \bibitem{c3}
        M. Santello, M. Flanders, and J. F. Soechting, ``Postural hand synergies for tool use,'' \emph{The Journal of Neuroscience}, vol. 18, no. 23, pp. 10105--10115, 1998.

        \bibitem{c4}
        A. d'Avella and E. Bizzi, ``Shared and specific muscle synergies in natural motor behaviors,'' \emph{Proceedings of the National Academy of Sciences}, vol. 102, no. 8, pp. 3076--3081, 2005.

        \bibitem{c5}
        M. Al Borno, J. L. Hicks, and S. L. Delp, ``The effects of motor modularity on performance, learning and generalizability in upper-extremity reaching: a computational analysis,'' \emph{Journal of The Royal Society Interface}, vol. 17, no. 167, Art. no. 20200011, 2020.

        \bibitem{c6}
        T. P. Lillicrap, J. J. Hunt, A. Pritzel, N. Heess, T. Erez, Y. Tassa, D. Silver, and D. Wierstra, ``Continuous control with deep reinforcement learning,'' in \emph{International Conference on Learning Representations (ICLR)}, 2016.

        \bibitem{c7}
        OpenAI, M. Andrychowicz, B. Baker, M. Chociej, R. Jozefowicz, B. McGrew, J. Pachocki, A. Petron, M. Plappert, G. Powell, \emph{et al.}, ``Learning dexterous in-hand manipulation,'' \emph{The International Journal of Robotics Research}, vol. 39, no. 1, pp. 3--32, 2020.

        \bibitem{c8}
        K. Kutsuzawa and M. Hayashibe, ``Motor synergy generalization framework for new targets in multi-planar and multi-directional reaching task,'' \emph{Royal Society Open Science}, vol. 9, no. 3, Art. no. 211721, 2022.

        \bibitem{c9}
        C. Berg, V. Caggiano, and V. Kumar, ``SAR: generalization of physiological agility and dexterity via synergistic action representation,'' \emph{Autonomous Robots}, vol. 48, no. 8, Art. no. 10182-4, 2024.

        \bibitem{c10}
        M. T. Ciocarlie and P. K. Allen, ``Hand posture subspaces for dexterous robotic grasping,'' \emph{The International Journal of Robotics Research}, vol. 28, no. 7, pp. 851--867, 2009.

        \bibitem{c11}
        M. G. Catalano, G. Grioli, E. Farnioli, A. Serio, C. Piazza, and A. Bicchi, ``Adaptive synergies for the design and control of the Pisa/IIT SoftHand,'' \emph{The International Journal of Robotics Research}, vol. 33, no. 5, pp. 768--782, 2014.

        \bibitem{c12}
        C. Della Santina, C. Piazza, G. Grioli, M. G. Catalano, and A. Bicchi, ``Toward dexterous manipulation with augmented adaptive synergies: The Pisa/IIT SoftHand 2,'' \emph{IEEE Transactions on Robotics}, vol. 34, no. 5, pp. 1141--1156, 2018.

        \bibitem{c13}
        J. Chai and M. Hayashibe, ``Motor synergy development in high-performing deep reinforcement learning algorithms,'' \emph{IEEE Robotics and Automation Letters}, vol. 5, no. 2, pp. 1271--1278, 2020.

        \bibitem{c14}
        J. Han, J. Chai, and M. Hayashibe, ``Synergy emergence in deep reinforcement learning for full-dimensional arm manipulation,'' \emph{IEEE Transactions on Medical Robotics and Bionics}, vol. 3, no. 2, pp. 498--509, 2021.

        \bibitem{c15}
        T. Haarnoja, A. Zhou, P. Abbeel, and S. Levine, ``Soft actor-critic: Off-policy maximum entropy deep reinforcement learning with a stochastic actor,'' in \emph{International Conference on Machine Learning (ICML)}, pp. 1861--1870, 2018.

        \bibitem{c16}
        M. Plappert, M. Andrychowicz, A. Ray, B. McGrew, B. Baker, G. Powell, J. Pachocki, W. Tworek, C. Cobbe, R. Vineet, \emph{et al.}, ``Multi-goal reinforcement learning: Challenging robotics environments and request for research,'' \emph{arXiv preprint arXiv:1802.09464}, 2018.

        \bibitem{c17}
        M. Andrychowicz, F. Wolski, A. Ray, J. Schneider, R. Fong, P. Welinder, B. McGrew, J. Tobin, P. Abbeel, and W. Zaremba, ``Hindsight experience replay,'' in \emph{Advances in Neural Information Processing Systems (NeurIPS)}, vol. 30, 2017.

        \bibitem{c18}
        R. de Lazcano, K. Andreas, J. J. Tai, S. R. Lee, and J. Terry, ``Gymnasium Robotics,'' in \emph{Farama Foundation}, 2024. 

        \bibitem{c19}
        A. Fukunishi, K. Kutsuzawa, D. Owaki, and M. Hayashibe, 
        ``Synergy quality assessment of muscle modules for determining 
        learning performance using a realistic musculoskeletal model,'' 
        \emph{Frontiers in Computational Neuroscience}, vol. 18, 
        Art. no. 1355855, 2024.

        \bibitem{c20} 
        T. Chen, J. Xu, and P. Agrawal, 
        ``A system for general in-hand object 
        re-orientation,'' 
        \emph{Proceedings of the 5th Conference on Robot Learning (CoRL)}, 2022, pp. 297-307.
                
        \bibitem{c21}
        E. Todorov and Z. Ghahramani, 
        ``Analysis of the synergies underlying 
        complex hand manipulation,''  
        \emph{Proceedings of the 26th Annual International Conference of the IEEE Engineering in Medicine and 
        Biology Society},
        vol. 2, 2004, pp. 4637-4640.
        
        \bibitem{c22}
        M. Berniker, A. Jarc, E. Bizzi, and M. C. Tresch, ``Simplified and 
        effective motor control based on muscle synergies to exploit 
        musculoskeletal dynamics,'' \emph{Proceedings of the National Academy 
        of Sciences}, vol. 106, pp. 7601--7606, 2009.

        \bibitem{c23}
        A. Rajeswaran, V. Kumar, A. Gupta, G. Vezzani, J. Schulman, E. Todorov, and 
        S. Levine, ``Learning complex dexterous manipulation with deep reinforcement 
        learning and demonstrations,'' in \emph{Robotics: Science and Systems (RSS)}, 
        2018.

\end{thebibliography}
        




\end{document}
